\documentclass[aspectratio=169]{beamer}

% \graphicspath{{figs}}  % your figures here

\usepackage{cmap}
\usepackage[T2A,T1]{fontenc}
\usepackage[utf8]{inputenc}
\usepackage[english,russian]{babel}
\usepackage{tikz}
\usetikzlibrary{positioning,arrows.meta}
\usepackage{array}
\usepackage{multicol}
\usepackage{booktabs}
\usepackage{csquotes}
\usepackage{amssymb,amsfonts,amsmath}

\usecolortheme{beaver}
\usefonttheme[onlymath]{serif}  % math fonts as in article
\setbeamertemplate{navigation symbols}{}
\setbeamertemplate{footline}[page number]
\setbeamertemplate{blocks}[rounded=true,shadow=true]
\setbeamersize{text margin left=0.5cm,text margin right=0.5cm}
\setbeamerfont{author}{size=\normalsize}
\setbeamerfont{institute}{size=\normalsize}
\setbeamerfont{date}{size=\normalsize}

% Обозначения, совпадающие с paper/main.tex и docs/problem_statement.md.
\newcommand{\I}{\mathbb{I}}
\newcommand{\Inv}{\mathrm{Inv}}
\newcommand{\Heval}{\mathcal{H}_{\mathrm{eval}}}
\newcommand{\fprod}{f_{\mathrm{prod}}}
\newcommand{\NA}{\mathrm{NA}}
\newcommand{\mcorr}{m_J^{\mathrm{corr}}}
\newcommand{\CP}{\mathrm{CP}}
\newcommand{\AR}{\mathrm{AR}}
\newcommand{\FAR}{\mathrm{FAR}}

\title{Система оценки качества ответов больших языковых моделей по новостным текстам стран Восточной Азии}
\subtitle{A System for Evaluating the Quality of Large Language Model Answers on East Asian News Texts}
\author{
    Багров Александр Михайлович\\
    Научный руководитель: Воронцов Константин Вячеславович, д.\,ф.-м.\,н.
}
\institute{
    Факультет вычислительной математики и кибернетики МГУ имени М.\,В.~Ломоносова\\
    Проект по курсу: оформление по шаблону <<Моя первая научная статья>>
}
\date{2026}

\begin{document}

\begin{frame}
    \thispagestyle{empty}
    \maketitle
\end{frame}

\begin{frame}{Система оценки качества ответов больших языковых моделей}
    Промышленная RAG-система <<Chinese Media Analysis>> (ИВ~РАН) отвечает по-русски на вопросы по 23,9~млн фрагментов китайскоязычных СМИ, но качество её ответов до сих пор оценивалось только качественно.
    \begin{block}{Цель}
        Построить систему оценки качества \emph{ответов} $\hat a=f(q)$ на вопросы о новостном корпусе: эталонную выборку, валидированные метрики и статистический протокол сравнения.
    \end{block}
    \begin{block}{Проблема}
        Готовые средства оценки RAG опираются на невалидированного LLM-судью; существующие наборы вопросов устаревают, так как ответы зависят от времени; сравнение промышленной системы с открытыми моделями не проводилось.
    \end{block}
    \begin{block}{Решение}
        \begin{enumerate}[1)]
            \item Ограничиться статическими вопросами: $\Inv(q)$~--- ответ не меняется при $t\ge t_0$.
            \item Построить $G=G^+\sqcup G^{\varnothing}$ с разметкой $R=3$ и согласием Криппендорфа $\alpha$.
            \item Отобрать метрики $M'\subset M$ по согласию с людьми и сравнить 13~систем по $Q(f)$.
        \end{enumerate}
    \end{block}
\end{frame}

\begin{frame}{Работа на одном слайде}
    От снимка корпуса к таблице результатов: шесть этапов, каждый воспроизводим по опубликованным конфигурациям.
    \begin{center}
    \begin{tikzpicture}[
        node distance=2.5mm,
        stage/.style={draw, rounded corners, align=center, font=\scriptsize,
                      text width=2.05cm, minimum height=1.15cm, inner sep=2pt},
        >={Stealth[length=2mm]}
    ]
        \node[stage] (a) {Снимок корпуса $D_{t_0}$, стратифицированный отбор};
        \node[stage, right=of a] (b) {Черновики $(q,a,S)$ моделью $g_{\mathrm{draft}}\notin\mathcal{L}$};
        \node[stage, right=of b] (c) {Разметка $R=3$, правила $E_1$--$E_5$, согласие $\alpha$};
        \node[stage, right=of c] (d) {Эталонная выборка $G=G^+\sqcup G^{\varnothing}$};
        \node[stage, right=of d] (e) {Метрики $M$, валидация $M'$ по $\bar h$};
        \node[stage, right=of e] (f) {Прогоны $\Heval$, статистика, $Q(f)$};
        \draw[->] (a) -- (b);
        \draw[->] (b) -- (c);
        \draw[->] (c) -- (d);
        \draw[->] (d) -- (e);
        \draw[->] (e) -- (f);
    \end{tikzpicture}
    \end{center}
    \begin{columns}
    \begin{column}{0.5\textwidth}
        \begin{enumerate}[1)]
            \item Обозначения: $D$, $\mathcal{Q}$, $\mathcal{A}=\mathcal{A}^+\sqcup\{\varnothing\}$, $\mathcal{I}$, $\Inv(q)$, $G$, $\mathcal{H}$, $M$, $Q(f)$.
            \item Объекты: ответ $\hat a=f(q)$~--- объект измерения; система $f$ входит только через $f(q)$.
        \end{enumerate}
    \end{column}
    \begin{column}{0.5\textwidth}
        \begin{enumerate}[1)]
            \item Данные: $n=200$--$300$ статических вопросов, $n_\varnothing\approx0{,}1\,n$.
            \item Системы: $\fprod$ и 6~открытых моделей в двух режимах ($|\Heval|=13$).
        \end{enumerate}
    \end{column}
    \end{columns}
    Следствие: получаем долговечный эталон, метрики с известным согласием с человеком и таблицу сравнения с интервалами и поправкой на множественные сравнения.
\end{frame}

\begin{frame}{Постановка задачи}
    \small
    \textbf{Дано.} Снимок корпуса $D=\{d_j\}_{j=1}^N$, $d_j=(\iota_j,u_j,\tau_j,s_j,x_j^{zh},x_j^{ru})$, $\tau_j\le t_0$; пространства $\mathcal{Q}$, $\mathcal{A}=\mathcal{A}^+\sqcup\{\varnothing\}$, $2^{\mathcal{I}}$;
    \[
        \Inv(q)\iff\exists\,a^*(q)\in\mathcal{A}^+:\ \forall t\ge t_0\quad a^*(q,t)=a^*(q),
        \qquad
        \widehat{\Inv}(q)=\I\Bigl[\textstyle\sum_{r=1}^{R} z_r(q)=R\Bigr].
    \]
    Эталонная выборка $G=G^+\sqcup G^{\varnothing}$, $G^+=\{(q_i,A^*_i,S^*_i,\mu_i)\}_{i=1}^{n_+}$, $\widehat{\Inv}(q_i)=1$.
    Класс гипотез $\mathcal{H}=\{f:\mathcal{Q}\to\mathcal{A}\times2^{\mathcal{I}}\}$,
    $\Heval=\{\fprod\}\cup\{f_0^{(m)}\}_{m\in\mathcal{L}}\cup\{f_\rho^{(m)}\}_{m\in\mathcal{L}}$.
    Метрики $M=\{m_k\}$, $m_k:\mathcal{A}\times2^{\mathcal{I}}\times G_i\to[0,1]\cup\{\NA\}$; оценки людей $\bar h_{i,f}$ на $G_h\subset G^+$.

    \medskip
    \textbf{Найти.} Выборку $G$; валидированное подсемейство метрик и лучшую систему:
    \[
        M'=\{m_k\in M:\ \rho_S(m_k,\bar h)\ge\rho_{\min}\},\ \rho_{\min}=0{,}5;\qquad
        Q(f)=\sum_{k} w_k\,\bar m_k(f),\qquad
        f^\star=\arg\max_{f\in\Heval} Q(f).
    \]

    \textbf{Критерий.} Первичная метрика $\mcorr$ (иначе F1); парный бутстреп по вопросам ($B=10^4$), критерий Уилкоксона, поправка Холма, уровень значимости $0{,}05$.

    \medskip
    \textbf{Основная гипотеза} $H_1$: $Q(\fprod)>Q(f_0^{(m)})$ для всех $m\in\mathcal{L}$.
\end{frame}

\begin{frame}{Решение: методика построения системы оценки}
    \begin{enumerate}[1)]
        \item \textbf{Отбор и черновики.} Стратифицированный отбор документов (домен $\times$ месяц $\times$ категория, переизбыток $\times2$); черновики $(q,a,S,\text{evidence})$ моделью $g_{\mathrm{draft}}\notin\mathcal{L}$; 388 старых вопросов~--- исходный пул.
        \item \textbf{Разметка.} $R=3$ независимых аннотатора (не менее одного читающего по-китайски); правила исключения $E_1$--$E_5$; классы FreshQA: never-changing и slow-changing с датой-якорем.
        \item \textbf{Согласие и сборка $G$.} Криппендорф $\alpha\ge0{,}67$; попарный F1 ответов $\ge0{,}8$; включение при единогласии, иначе адъюдикация; пилот 30 $\to$ полный прогон $\to$ заморозка \texttt{v1.0}.
        \item \textbf{Метрики $M$.} По эталону: EM, F1, BERTScore, $\mcorr$ (новое измерение); цитирование: $\CP_{\mathrm{id}}$, $\mathrm{CR}_{\mathrm{id}}$, $\CP_{\mathrm{ent}}$; отказы: $\AR$, $1-\FAR$; поиск: Hit@$k$, Recall@$k$, MRR, nDCG@$k$. Существующий семимерный судья~--- кандидатное подсемейство.
        \item \textbf{Валидация $M'$.} Оценки людей на $G_h$ ($n_h=100$); отбор по $\rho_S\ge0{,}5$ с бутстреп-интервалом; проверка самопредпочтения судьи ($H_5$).
        \item \textbf{Прогоны и статистика.} 13 систем при temperature~$0$; $Q(f)$, парный бутстреп, Уилкоксон, Холм, $\delta$ Клиффа; разбивка по стратам.
    \end{enumerate}
\end{frame}

\begin{frame}{План экспериментов}
    \begin{columns}
        \begin{column}{0.42\textwidth}
            \small
            \textbf{Системы} ($|\Heval|=13$): $\fprod$ (Qdrant + SearchAPI + генератор) и $\mathcal{L}$~--- шесть открытых моделей провайдера ИВ~РАН в режимах без поиска $f_0^{(m)}$ и с общим ретривером $f_\rho^{(m)}$.

            \medskip
            \textbf{Размеры:} $n=200$--$300$ вопросов ($n\ge150$), $n_\varnothing\approx0{,}1\,n$, $R=3$, $n_h=100$.

            \medskip
            \textbf{Бюджет:} разметка $\approx21$~ч на аннотатора; $\approx4750$ генераций ($\approx40$~ч); судья $\approx25$~ч.

            \medskip
            \textbf{Гипотезы:} $H_1$--$H_3$~--- сравнение по $Q(f)$; $H_4$~--- $\rho_S(\mcorr,\bar h)\ge0{,}5$, $\alpha\ge0{,}67$; $H_5$~--- самопредпочтение судьи незначимо.
        \end{column}
        \begin{column}{0.56\textwidth}
            \centering
            \scriptsize
            \begin{tabular}{@{}llccccc@{}}
                \toprule
                Система & Режим & $\mcorr$ & F1 & $\CP_{\mathrm{id}}$ & $\AR$ & $Q(f)$ \\
                \midrule
                $\fprod$ & Qdrant + SearchAPI & $\dots$ & $\dots$ & $\dots$ & $\dots$ & $\dots$ \\
                \midrule
                $f_0^{(m_1)}$ & без поиска & $\dots$ & $\dots$ & $\NA$ & $\dots$ & $\dots$ \\
                $\vdots$ & $\vdots$ & $\vdots$ & $\vdots$ & $\vdots$ & $\vdots$ & $\vdots$ \\
                $f_0^{(m_6)}$ & без поиска & $\dots$ & $\dots$ & $\NA$ & $\dots$ & $\dots$ \\
                \midrule
                $f_\rho^{(m_1)}$ & $\rho_{\mathrm{prod}}$ & $\dots$ & $\dots$ & $\dots$ & $\dots$ & $\dots$ \\
                $\vdots$ & $\vdots$ & $\vdots$ & $\vdots$ & $\vdots$ & $\vdots$ & $\vdots$ \\
                $f_\rho^{(m_6)}$ & $\rho_{\mathrm{prod}}$ & $\dots$ & $\dots$ & $\dots$ & $\dots$ & $\dots$ \\
                \bottomrule
            \end{tabular}

            \medskip
            \tiny В ячейках~--- $\bar m_k(f)$ с 95\,\%-м бутстреп-интервалом; звёздочка~--- значимое отличие от $\fprod$ после поправки Холма. Значения появятся после прогонов (январь--февраль 2027).
        \end{column}
    \end{columns}
\end{frame}

\begin{frame}{Заключение}
    \begin{enumerate}[1)]
        \item Будет формализована задача оценки качества ответов на статические вопросы по новостному корпусу: предикат $\Inv(q)$, класс гипотез $\Heval$, семейство метрик $M$, показатель $Q(f)$ и критерий сравнения.
        \item Будет построена и опубликована эталонная выборка $G=G^+\sqcup G^{\varnothing}$ ($n=200$--$300$) с разметкой $R=3$ аннотаторами, мерой согласия $\alpha$ и стабильными идентификаторами документов.
        \item Будет получено валидированное подсемейство метрик $M'$, включающее метрику правильности по эталону $\mcorr$ и метрики цитирования и отказов, с оценкой согласия с человеком $\rho_S$.
        \item Будет проведено первое систематическое сравнение промышленной системы с открытыми моделями (13~систем) с бутстреп-интервалами и поправкой Холма; проверены гипотезы $H_1$--$H_5$.
    \end{enumerate}
\end{frame}

\end{document}
